% theme: observation_and_classification
\MajorQuestion{生物の観察／生物の特徴と分類のしかた}
\SetRowSpaceQanda

\Instruction{次の問いに答えなさい。}
\QQ{両目で立体的に観察できる、20～40倍程度の器具を何というか。}{双眼実体顕微鏡}
\QQ{野外で生物を5～10倍程度に拡大して見る器具を\NamedBlank{loupe}という。}{ルーペ}
\QQ{スライドガラスに観察物をのせ、カバーガラスをかけたものを何というか。}{プレパラート}
\QQ{40～600倍程度まで拡大して観察できる器具を\NamedBlank{microscope}という。}{顕微鏡}
\QQ{生物を大きさや形などの特徴を基準にして仲間分けすることを何というか。}{分類}

\Instruction{\strut}
\QQ{\RefBlank{loupe}で観察すると目を痛めるため、見てはいけないものは何か。}{太陽}
\QQ{\RefBlank{microscope}で接眼レンズが10倍、対物レンズが40倍のとき、全体の倍率は何倍か。}{400倍}
\QQ{観察したもののスケッチは、点とどのような線だけではっきりとかくか。}{細い線}
\QQ{観察物を動かせるとき、\RefBlank{loupe}を目に近づけたまま何を動かすか。}{観察するもの}
\QQ{\RefBlank{microscope}の視野では、観察物は実物と比べてどの向きに見えるか。}{上下左右が逆向き}
